\documentclass[11pt]{article}

\usepackage{amsmath,amsthm,amssymb}
\usepackage[margin=1in]{geometry}

\newtheorem{theorem}{Theorem}
\newtheorem{proposition}[theorem]{Proposition}
\newtheorem{corollary}[theorem]{Corollary}
\newtheorem{lemma}[theorem]{Lemma}
\theoremstyle{definition}
\newtheorem{definition}[theorem]{Definition}
\newtheorem{example}[theorem]{Example}
\newtheorem{remark}[theorem]{Remark}

\newcommand{\Hom}{\operatorname{Hom}}
\newcommand{\End}{\operatorname{End}}
\newcommand{\id}{\mathrm{id}}
\newcommand{\C}{\mathcal{C}}
\newcommand{\D}{\mathcal{D}}
\newcommand{\ob}{\operatorname{ob}}
\newcommand{\Set}{\mathbf{Set}}

\title{Collage reconstruction: a tightened proof\\
\large Bisets over endomorphism monoids, the tensor coequalizer,\\
and the reconstruction of a small category}
\author{MacBeth\\\small Kodamai}
\date{10 June 2026}

\begin{document}
\maketitle

\begin{abstract}
This note replaces the hand-waved proof of the ``collage reconstruction''
theorem (part (1) of the Corrected Decomposition theorem in the two-atoms /
Zappa--Sz\'ep note) by a fully element-wise argument. We define the tensor
product of bisets over a monoid as an explicit coequalizer, prove that
composition in a small category is balanced and therefore descends to the
tensor, verify the category axioms on representatives, and prove that the
collage construction and its converse are mutually inverse. One genuine
subtlety in the converse --- the need for the regular-biset embedding of each
vertex monoid to be \emph{strictly} compatible with composition --- is
isolated and discharged explicitly rather than glossed.
\end{abstract}

\section{The tensor product of bisets over a monoid}

Throughout, $M,N,P$ are monoids (written multiplicatively, units $1$).
A \emph{right $M$-set} is a set $X$ with a map $X\times M\to X$,
$(x,m)\mapsto x\cdot m$, satisfying $x\cdot 1 = x$ and
$(x\cdot m)\cdot m' = x\cdot(mm')$. Left $M$-sets are dual. An
\emph{$(N,M)$-biset} is a set carrying a left $N$-action and a right
$M$-action that commute: $(n\cdot x)\cdot m = n\cdot(x\cdot m)$.

\begin{definition}[Tensor over a monoid as a coequalizer]\label{def:tensor}
Let $X$ be a right $M$-set and $Y$ a left $M$-set. Consider the two maps
\[
  \lambda,\rho : X\times M\times Y \;\rightrightarrows\; X\times Y,
  \qquad
  \lambda(x,m,y) = (x\cdot m,\, y),
  \quad
  \rho(x,m,y) = (x,\, m\cdot y).
\]
The \emph{tensor product} $X\otimes_M Y$ is the coequalizer of
$\lambda,\rho$ in $\Set$: the quotient $(X\times Y)/{\sim}$, where $\sim$ is
the smallest equivalence relation with
\[
  (x\cdot m,\, y) \;\sim\; (x,\, m\cdot y)
  \qquad\text{for all } x\in X,\; m\in M,\; y\in Y .
\]
We write $x\otimes y$ for the class of $(x,y)$ and
$q:X\times Y\to X\otimes_M Y$ for the quotient map. By construction
$x\cdot m\otimes y = x\otimes m\cdot y$.
\end{definition}

\begin{remark}\label{rem:relation}
The generating relation is already reflexive ($m=1$), but it is neither
symmetric nor transitive in general, so $\sim$ is its
reflexive--symmetric--transitive closure: $(x,y)\sim(x',y')$ iff there is a
finite zig-zag
\[
  (x,y) = (x_0,y_0) \,\leftrightarrow\, (x_1,y_1)
  \,\leftrightarrow\, \cdots \,\leftrightarrow\, (x_k,y_k) = (x',y'),
\]
each step $\leftrightarrow$ being a generator applied left-to-right or
right-to-left. We will need this explicit zig-zag description when checking
that maps out of the tensor are well defined.
\end{remark}

The universal property is the defining one for a coequalizer; we record the
form we use.

\begin{lemma}[Universal property / balance]\label{lem:up}
Let $X$ be a right $M$-set, $Y$ a left $M$-set, $Z$ any set. A map
$\beta:X\times Y\to Z$ \emph{descends} to a (necessarily unique) map
$\bar\beta:X\otimes_M Y\to Z$ with $\bar\beta(x\otimes y)=\beta(x,y)$ if and
only if $\beta$ is \emph{balanced}:
\[
  \beta(x\cdot m,\, y) = \beta(x,\, m\cdot y)
  \qquad\text{for all } x\in X,\ m\in M,\ y\in Y. \tag{B}
\]
\end{lemma}

\begin{proof}
($\Rightarrow$) If $\bar\beta$ exists then
$\beta(x\cdot m,y)=\bar\beta(x\cdot m\otimes y)
=\bar\beta(x\otimes m\cdot y)=\beta(x,m\cdot y)$.

($\Leftarrow$) Suppose (B). Then $\beta$ is constant on each generator pair,
hence (being a function landing in a set, so respecting symmetry and
composing along the zig-zag of Remark~\ref{rem:relation}) constant on each
$\sim$-class: if $(x,y)\leftrightarrow(x',y')$ is a single generator step in
either direction, (B) gives $\beta(x,y)=\beta(x',y')$, and chaining over the
finite zig-zag gives equality on the whole class. Therefore
$\bar\beta([x,y]):=\beta(x,y)$ is well defined; it satisfies
$\bar\beta\circ q=\beta$, and uniqueness is forced because $q$ is surjective.
\end{proof}

\section{The collage of a small category}

Fix a small category $\C$. For objects $a,b$ write $\Hom(a,b)$ for the
hom-set and $\End(a):=\Hom(a,a)$ for the endomorphism monoid (composition,
unit $\id_a$). As in the source note, $\Hom(a,b)$ is canonically an
$(\End(b),\End(a))$-biset: the right $\End(a)$-action is precomposition
$f\cdot e := f\circ e$ and the left $\End(b)$-action is postcomposition
$e'\cdot f := e'\circ f$; the two commute by associativity, so the biset
axiom holds.

\subsection*{(i) Balance}

\begin{proposition}[Composition is balanced]\label{prop:balance}
Fix objects $a,b,c$. The composition map
\[
  \gamma : \Hom(b,c)\times\Hom(a,b)\to\Hom(a,c),
  \qquad \gamma(g,f) = g\circ f,
\]
is $\End(b)$-balanced: for all $g\in\Hom(b,c)$, $e\in\End(b)$,
$f\in\Hom(a,b)$,
\[
  \gamma(g\cdot e,\, f) = \gamma(g,\, e\cdot f),
  \qquad\text{i.e.}\qquad (g\cdot e)\circ f = g\circ(e\cdot f).
\]
Hence by Lemma~\ref{lem:up} it descends to a unique map
\[
  \mu_{a,b,c} : \Hom(b,c)\otimes_{\End(b)}\Hom(a,b)\longrightarrow\Hom(a,c),
  \qquad \mu_{a,b,c}(g\otimes f) = g\circ f .
\]
\end{proposition}

\begin{proof}
Here $g\cdot e = g\circ e$ (right $\End(b)$-action on $\Hom(b,c)$) and
$e\cdot f = e\circ f$ (left $\End(b)$-action on $\Hom(a,b)$). Associativity of
composition in $\C$ gives
\[
  (g\cdot e)\circ f = (g\circ e)\circ f = g\circ(e\circ f) = g\circ(e\cdot f),
\]
which is exactly condition (B) of Lemma~\ref{lem:up} for $\beta=\gamma$ and
$M=\End(b)$. The descended map $\mu_{a,b,c}$ and its formula are then immediate
from the lemma.
\end{proof}

\subsection*{(ii) Reconstruction}

\begin{definition}[The collage category $\C^{\flat}$]\label{def:collage}
Define a structure $\C^{\flat}$ with
\begin{itemize}
\item objects: $\ob\C$;
\item hom-sets: $\C^{\flat}(a,b) := \Hom(a,b)$ (the underlying set of the
biset);
\item composition: $\mu_{a,b,c}$ of Proposition~\ref{prop:balance}, precomposed
with the quotient $q$, i.e. $g\circ^{\flat} f := \mu_{a,b,c}(g\otimes f)$;
\item identities: $\id^{\flat}_a := 1_{\End(a)} = \id_a$, the monoid unit.
\end{itemize}
\end{definition}

\begin{theorem}[Reconstruction]\label{thm:reconstruct}
$\C^{\flat}$ is a category, and the object-identity assignment
$F:\C\to\C^{\flat}$, $F(a)=a$ on objects and $F(f)=f$ on morphisms, is an
isomorphism of categories with inverse the identity-on-everything map.
\end{theorem}

\begin{proof}
We verify the axioms element-wise; since
$\C^{\flat}(a,b)=\Hom(a,b)$ as sets and, by the formula in
Proposition~\ref{prop:balance},
\[
  g\circ^{\flat} f \;=\; \mu_{a,b,c}(g\otimes f) \;=\; g\circ f ,
\]
\emph{composition in $\C^{\flat}$ is literally composition in $\C$}, and
$\id^{\flat}_a=\id_a$. Thus every axiom of $\C^{\flat}$ is the corresponding
axiom of $\C$:
\begin{description}
\item[Well-definedness.] $\mu_{a,b,c}$ is a genuine function on the tensor by
Proposition~\ref{prop:balance}, and $q$ is a function, so $\circ^{\flat}$ is a
well-defined map $\Hom(b,c)\times\Hom(a,b)\to\Hom(a,c)$.
\item[Left unit.] $\id^{\flat}_b\circ^{\flat} f = \id_b\circ f = f$ for
$f\in\Hom(a,b)$.
\item[Right unit.] $f\circ^{\flat}\id^{\flat}_a = f\circ\id_a = f$.
\item[Associativity.] For $f\in\Hom(a,b)$, $g\in\Hom(b,c)$,
$h\in\Hom(c,d)$,
\[
  (h\circ^{\flat} g)\circ^{\flat} f = (h\circ g)\circ f
  = h\circ(g\circ f) = h\circ^{\flat}(g\circ^{\flat} f),
\]
by associativity in $\C$.
\end{description}
Hence $\C^{\flat}$ is a category. The map $F$ is the identity on objects and on
each hom-set, and it preserves composition and identities because
$g\circ^{\flat} f = g\circ f$ and $\id^{\flat}_a=\id_a$. A bijective-on-objects,
bijective-on-homs functor is an isomorphism of categories, with the
set-theoretic identity as inverse functor.
\end{proof}

\begin{remark}[Where the tensor actually does work]
The tensor $\otimes_{\End(b)}$ is not cosmetic: it is exactly the data needed
to phrase composition without reference to the ambient category, by quotienting
the formal pair-set $\Hom(b,c)\times\Hom(a,b)$ by the $\End(b)$-loop relation.
The non-trivial content of Proposition~\ref{prop:balance} is that the relation
imposed by the coequalizer is \emph{respected} by composition; once that holds,
$\mu$ exists and the category axioms transfer verbatim. In particular the often
mysterious ``associativity is compatibility of the tensor glueing'' reduces, on
representatives, to plain associativity in $\C$ --- because $\mu(g\otimes f)$ is
computed by a chosen representative and the value $g\circ f$ is
representative-independent by balance.
\end{remark}

\subsection*{(iii) Converse: from a collage datum to a category}

We now axiomatize the data abstractly and show the construction inverts.

\begin{definition}[Biset-collage datum]\label{def:datum}
A \emph{biset-collage datum} $\D$ consists of:
\begin{enumerate}
\item a set $O$ of objects;
\item for each $x\in O$ a monoid $M_x$ (the vertex monoid), with unit $1_x$;
\item for each ordered pair $(a,b)$ an $(M_b,M_a)$-biset $H(a,b)$, with a
\emph{distinguished} biset isomorphism $\iota_x:M_x \xrightarrow{\ \sim\ }
H(x,x)$ for $a=b=x$, where $M_x$ is viewed as the regular $(M_x,M_x)$-biset
(left and right action $=$ multiplication); we write
$\underline e := \iota_x(e)\in H(x,x)$ and identify $M_x$ with its image, so
that $\underline{1_x}\in H(x,x)$ is a distinguished element and the biset
actions of $M_x$ on every $H(a,b)$ agree with multiplication on $H(x,x)$ in
the sense of (U) below;
\item for each triple $(a,b,c)$ a \emph{balanced} composition
\[
  c_{a,b,c} : H(b,c)\times H(a,b)\to H(a,c),
  \qquad c(g,e\cdot f)=c(g\cdot e, f)\ \ (e\in M_b),
\]
inducing $m_{a,b,c}:H(b,c)\otimes_{M_b}H(a,b)\to H(a,c)$ by
Lemma~\ref{lem:up};
\end{enumerate}
subject to the axioms (writing $g\bullet f := c(g,f)$):
\begin{description}
\item[(U) Unit.] For $f\in H(a,b)$,
$\underline{1_b}\bullet f = f = f\bullet\underline{1_a}$, and moreover the
distinguished left/right actions are recovered by composition with elements of
the vertex monoids: for $e\in M_b$, $\underline e\bullet f = e\cdot f$, and for
$e\in M_a$, $f\bullet\underline e = f\cdot e$. (Equivalently: on
$H(b,b)\times H(a,b)$, $\bullet$ restricts to the left $M_b$-action under
$\iota_b$, and dually on the right.)
\item[(A) Associativity.] For $f\in H(a,b)$, $g\in H(b,c)$, $h\in H(c,d)$,
$(h\bullet g)\bullet f = h\bullet(g\bullet f)$.
\end{description}
\end{definition}

\begin{theorem}[Converse and mutual inverse]\label{thm:converse}
\leavevmode
\begin{enumerate}
\item \emph{(Reconstruction.)} Every biset-collage datum $\D$ yields a small
category $\C_\D$ with $\ob\C_\D=O$, $\C_\D(a,b)=H(a,b)$, composition
$g\circ f := g\bullet f$, and identities $\id_a := \underline{1_a}$.
\item \emph{(Mutual inverse.)} The two constructions are inverse up to
isomorphism: $(\C^{\flat})$-of-$\D$ on the collage datum extracted from a
category returns that category (Theorem~\ref{thm:reconstruct}), and conversely
$\C_{\D(\C_\D)} = \C_\D$ on the nose, where $\D(\C_\D)$ is the collage datum
extracted from $\C_\D$ by the recipe of Section~2 ($M_x=\End_{\C_\D}(x)$,
$H=\Hom_{\C_\D}$, $\iota_x=\id$).
\end{enumerate}
\end{theorem}

\begin{proof}
(1) Axioms (U) and (A) are precisely the unit and associativity laws of a
category, stated for $\bullet$, with $\underline{1_a}$ as identity at $a$.
Well-definedness of composition as a map $H(b,c)\times H(a,b)\to H(a,c)$ is part
of the datum (item 4); balance is not even needed for $\C_\D$ to be a category,
but it guarantees that $\bullet$ factors through
$H(b,c)\otimes_{M_b}H(a,b)$, which is what makes $\C_\D$ a genuine collage
rather than an arbitrary composition law. Smallness holds because $O$ and all
$H(a,b)$ are sets.

\medskip
\noindent(2) \emph{Category $\to$ datum $\to$ category.} Starting from a small
category $\C$, the extracted datum has $O=\ob\C$, $M_x=\End(x)$,
$H(a,b)=\Hom(a,b)$ with the canonical biset structure, $\iota_x=\id$ (the
regular biset of $\End(x)$ \emph{is} $\Hom(x,x)$ with its pre/post-composition
actions), and $c_{a,b,c}=\gamma$ (composition), which is balanced by
Proposition~\ref{prop:balance}. Axiom (U) for this datum says
$\id_b\circ f=f=f\circ\id_a$ and $\underline e\bullet f=e\circ f=e\cdot f$:
the first is the unit law of $\C$, the second is the definition of the left
action; both hold. Axiom (A) is associativity in $\C$. Thus the extracted
datum is a legitimate biset-collage datum, and $\C_{\D(\C)}$ has the same
objects, the same hom-sets, composition $g\bullet f=g\circ f$ and identities
$\underline{1_a}=\id_a$ --- i.e. $\C_{\D(\C)}=\C$ on the nose. (This is the
content of Theorem~\ref{thm:reconstruct} read backwards.)

\medskip
\noindent\emph{Datum $\to$ category $\to$ datum.} Starting from a datum $\D$,
form $\C_\D$ as in (1); then extract a datum $\D':=\D(\C_\D)$. We check
$\D'=\D$. Objects: $\ob\C_\D=O$. Vertex monoids: $\End_{\C_\D}(x)=H(x,x)$ as a
\emph{set}, with multiplication $g\bullet f$; by axiom (U) restricted to
$x=a=b=c$, for $\underline e,\underline{e'}\in\iota_x(M_x)$ we have
$\underline{e'}\bullet\underline e=e'\cdot\underline e=\underline{e'e}$ (left
action $=$ multiplication via $\iota_x$, since $\iota_x$ is a biset iso onto the
regular biset). Because $\iota_x:M_x\to H(x,x)$ is a bijection, every element of
$H(x,x)$ is some $\underline e$, so multiplication in $\End_{\C_\D}(x)$ is
$\underline{e'}\bullet\underline e=\underline{e'e}$, i.e.
$\iota_x$ is a monoid isomorphism $M_x\cong\End_{\C_\D}(x)$; identifying along
$\iota_x$ (as the datum already does), $M_x=\End_{\C_\D}(x)$ with equal units
$\underline{1_x}=\id_x$. Hom-sets: $H(a,b)=\Hom_{\C_\D}(a,b)$ by definition.
Biset actions: the extracted left $\End_{\C_\D}(b)$-action on $\Hom_{\C_\D}(a,b)$
is postcomposition $\underline e\bullet f$, which by axiom (U) equals the
original left action $e\cdot f$; dually on the right. Hence the extracted biset
structure equals the original. Composition: $\bullet$ is unchanged. Therefore
$\D'=\D$, as required.
\end{proof}

\begin{remark}[The one genuine subtlety, isolated]\label{rem:subtle}
The single place where strictness could fail in the converse is the clause
``$\iota_x$ is a biset isomorphism onto the \emph{regular} biset'' together with
axiom (U)'s identity $\underline e\bullet f=e\cdot f$. Without it, the
distinguished elements $\underline e\in H(x,x)$ would compose by some abstract
$\bullet$ that need not coincide with the monoid multiplication of $M_x$, and
then ``$\End_{\C_\D}(x)=M_x$'' would hold only up to a chosen isomorphism, not
on the nose. We have made this compatibility \emph{part of the datum}
(items 3--4 and axiom (U)); under it the round trip is an equality. If one
weakens the datum to allow $\iota_x$ to be merely a monoid \emph{homomorphism}
that is bijective on the underlying set but \emph{not} required to intertwine
$\bullet$ with $M_x$-multiplication, the two constructions are inverse only up
to canonical isomorphism of categories, not equal. We flag this explicitly: the
``mutually inverse on the nose'' statement is true exactly for data satisfying
the regular-embedding axiom (U), and is the natural normal form; for general
data one gets an equivalence (indeed isomorphism) of categories rather than an
identity. No gap remains in the stated (strict) version.
\end{remark}

\section{Summary of the logical skeleton}

\begin{enumerate}
\item \textbf{Tensor} $=$ coequalizer of the two action maps
(Definition~\ref{def:tensor}); maps out of it $=$ balanced maps
(Lemma~\ref{lem:up}).
\item \textbf{Balance} of composition $=$ associativity in $\C$
(Proposition~\ref{prop:balance}); hence $\mu(g\otimes f)=g\circ f$ is
well defined.
\item \textbf{Reconstruction:} $\C^{\flat}$ has composition literally equal to
that of $\C$, so the category axioms transfer verbatim and
$\C\cong\C^{\flat}$ (Theorem~\ref{thm:reconstruct}).
\item \textbf{Converse:} a biset-collage datum with the regular-embedding
axiom (U) and associativity (A) is exactly a small category, and the two
constructions are mutually inverse on the nose
(Theorem~\ref{thm:converse}), the only subtlety being the strict
compatibility of the vertex-monoid embedding, which is isolated in
Remark~\ref{rem:subtle}.
\end{enumerate}

\end{document}
